\documentclass[12pt,a4paper]{article}

\usepackage[utf8]{inputenc}
\usepackage[T1]{fontenc}
\usepackage[brazil]{babel}
\usepackage{geometry}
\usepackage{amsmath}
\usepackage{amssymb}
\usepackage{graphicx}
\usepackage{booktabs}
\usepackage{array}
\usepackage{hyperref}
\usepackage{xcolor}
\usepackage{listings}
\usepackage{caption}
\usepackage{float}
\usepackage{enumitem}
\usepackage{parskip}
\usepackage{tabularx}

\geometry{top=2.5cm, bottom=2.5cm, left=3cm, right=2.5cm}

\hypersetup{
    colorlinks=true,
    linkcolor=black,
    urlcolor=blue,
    citecolor=black
}

\lstset{
    basicstyle=\ttfamily\small,
    backgroundcolor=\color{gray!8},
    frame=single,
    framerule=0.4pt,
    rulecolor=\color{gray!40},
    breaklines=true,
    columns=fullflexible,
    keepspaces=true,
}

% ─── Dados do projeto ────────────────────────────────────────────────────────
\title{
    \vspace{-1cm}
    \large Teoria dos Grafos --- Relatório de Entrega\\[0.4em]
    \LARGE \textbf{Parte 1: Rede de Aeroportos Brasileiros}
}
\author{Anderson Gabriel \and Luís Filipe \and Rafael Peixoto \and Raphael Vasconcelos \and Yuri Cavalcanti}
\date{Junho de 2026}

\begin{document}

\maketitle
\tableofcontents
\newpage

% ─── 1. Descrição do Domínio ─────────────────────────────────────────────────
\section{Descrição do Domínio}

Esta entrega modela a malha aérea nacional brasileira como um grafo, utilizando dados reais de voos
realizados em janeiro de 2026, fornecidos pela ANAC (Agência Nacional de Aviação Civil) por meio do
dataset VRA\footnote{Voos Regulares da Aviação --- disponível no portal de dados abertos da ANAC.}.

O objetivo é representar os aeroportos como nós e as rotas operadas entre eles como arestas, permitindo
a aplicação de algoritmos de grafos para análise estrutural da malha (BFS, DFS) e busca de caminhos
mínimos (Dijkstra).

% ─── 2. Construção do Grafo ───────────────────────────────────────────────────
\section{Construção do Grafo}

\subsection{Nós --- Aeroportos}

Após limpeza e filtragem do dataset (voos do tipo \texttt{N} --- nacional, situação \texttt{REALIZADO}),
identificamos \textbf{144 aeroportos nacionais} únicos. Cada aeroporto é identificado pelo seu código
\textbf{ICAO} (ex.: \texttt{SBGR}, \texttt{SBEG}), utilizando a biblioteca Python
\texttt{airportsdata}\footnote{\url{https://pypi.org/project/airportsdata/} --- base de dados de
aeroportos mundiais, indexada por código ICAO, que fornece nome da cidade, estado, país e
coordenadas geográficas (latitude/longitude) de cada aeroporto.} como fonte de metadados.
A partir do código ICAO, a biblioteca resolve automaticamente cidade e estado de cada aeroporto,
permitindo o mapeamento para as cinco regiões geográficas brasileiras: Norte, Nordeste, Sudeste,
Sul e Centro-Oeste. Casos não resolvidos automaticamente (aeroportos com dados incompletos na
base) foram corrigidos manualmente no pipeline.

\subsection{Arestas --- Rotas e Direcionamento}

O grafo é \textbf{dirigido}. Essa escolha reflete a natureza dos dados: o dataset da ANAC registra
cada voo com origem e destino distintos como eventos independentes. Um voo \texttt{SBGR→SBBR} e um voo
\texttt{SBBR→SBGR} podem ter frequências, empresas e demandas diferentes --- colapsá-los em uma aresta
não-dirigida eliminaria essa assimetria real.

Exemplos que justificam o direcionamento:
\begin{itemize}
    \item Um aeroporto hub pode receber voos de 30 destinos, mas operar saídas para apenas 20 --- grau
          de entrada $\neq$ grau de saída.
    \item A frequência \texttt{BSB→BEL} pode diferir de \texttt{BEL→BSB} por demanda assimétrica
          (turismo, negócios, sazonalidade).
    \item Dijkstra opera naturalmente sobre grafos dirigidos, preservando a semântica
          correta de ``menor caminho de $A$ para $B$''.
\end{itemize}

O grafo resultante possui \textbf{955 arestas dirigidas}, cada uma representando um par ordenado
(origem, destino) com ao menos uma operação realizada em janeiro de 2026.

% ─── 3. Peso das Arestas ─────────────────────────────────────────────────────
\section{Peso das Arestas}

\subsection{Motivação: por que frequência de voos?}

A escolha da métrica de peso é a principal decisão de modelagem desta entrega. As alternativas
consideradas foram:

\begin{table}[H]
\centering
\caption{Alternativas de peso avaliadas}
\begin{tabularx}{\textwidth}{lXX}
\toprule
\textbf{Métrica} & \textbf{Vantagem} & \textbf{Problema} \\
\midrule
Distância geográfica (km) & Intuitiva, estável & Ignora conectividade operacional real \\[4pt]
Tempo médio de voo (h)    & Reflete duração real & Dados ausentes no dataset; correlaciona com km \\[4pt]
Número de empresas        & Indica concorrência & Granularidade insuficiente \\[4pt]
\textbf{Frequência de voos} & \textbf{Reflete demanda e oferta reais} & \textbf{---} \\
\bottomrule
\end{tabularx}
\end{table}

A \textbf{frequência de voos} (número de operações realizadas no período) foi escolhida porque:

\begin{enumerate}
    \item É diretamente extraível do dataset da ANAC sem dependências externas.
    \item Captura simultaneamente oferta (quantas empresas operam) e demanda (quantas decolagens
          aconteceram de fato).
    \item Pares de alta frequência correspondem a conexões operacionalmente ``mais próximas'' ---
          mais voos significa mais opções de embarque, menor espera efetiva e maior confiabilidade
          da rota.
    \item A distância geográfica e o tempo de voo são altamente correlacionados entre si e omitem
          a dimensão de conectividade da malha. Dois aeroportos fisicamente próximos podem ter
          zero voos entre si; dois distantes podem ter 50 voos diários.
\end{enumerate}

\subsection{Fórmula do Peso}

O peso deve representar \textbf{custo de travessia}: quanto maior, mais ``caro'' é usar aquela
conexão. Como frequência alta indica conexão melhor, o peso é o \emph{inverso} da frequência.

Para evitar escala linear excessiva (um par com 1000 voos não deve ter peso 1000 vezes menor que
um par com 1 voo), aplicamos o logaritmo natural antes de inverter:

\begin{equation}
    w(u, v) = \frac{1}{\ln(f_{uv} + 1)}
    \label{eq:peso}
\end{equation}

onde $f_{uv}$ é o número de voos realizados do aeroporto $u$ para o aeroporto $v$ em janeiro de 2026,
e o $+1$ evita $\ln(0)$ no caso de frequência zero.

\subsubsection{Implementação}

Em Python, utilizamos \texttt{math.log1p}, que computa $\ln(x+1)$ com maior precisão numérica para
valores próximos de zero:

\begin{lstlisting}[language=Python]
from math import log1p

edge_df['peso'] = edge_df['quantidade'].apply(
    lambda freq: round(1 / log1p(freq), 4)
)
\end{lstlisting}

\subsubsection{Comportamento da fórmula}

A Tabela~\ref{tab:pesos} ilustra como o peso varia com a frequência:

\begin{table}[H]
\centering
\caption{Relação entre frequência de voos e peso da aresta}
\label{tab:pesos}
\begin{tabular}{rrl}
\toprule
\textbf{Frequência ($f$)} & \textbf{$w = 1/\ln(f+1)$} & \textbf{Interpretação} \\
\midrule
1   & 1{,}4427 & conexão rara \\
5   & 0{,}5581 & conexão ocasional \\
9   & 0{,}4343 & conexão regular \\
13  & 0{,}3789 & conexão frequente \\
31  & 0{,}2885 & hub regional \\
48  & 0{,}2569 & hub principal \\
100 & 0{,}2171 & corredor de alta demanda \\
\bottomrule
\end{tabular}
\end{table}

\subsubsection{Exemplos reais do dataset}

\begin{table}[H]
\centering
\caption{Arestas reais com peso calculado}
\begin{tabular}{llcrl}
\toprule
\textbf{Origem} & \textbf{Destino} & \textbf{Tipo} & \textbf{Voos} & \textbf{Peso} \\
\midrule
SBAE & SBCF & intrarregional & 1  & 1{,}4427 \\
SBAE & SBMO & interregional  & 9  & 0{,}4343 \\
SBAE & SBKP & intrarregional & 31 & 0{,}2885 \\
SBAR & SBGL & interregional  & 48 & 0{,}2569 \\
SBAR & SBBR & interregional  & 35 & 0{,}2791 \\
\bottomrule
\end{tabular}
\end{table}

O logaritmo comprime a escala: a diferença de peso entre 1 voo (1{,}4427) e 48 voos (0{,}2569)
é de apenas $\approx 1{,}19$ unidades, ao invés de 47 unidades sem o log. Isso impede que conexões
com altíssima frequência dominem completamente os caminhos mínimos a ponto de tornar irrelevantes
todas as outras arestas.

\subsection{Densidade do Grafo}

O grafo é dirigido, portanto a fórmula correta de densidade considera todos os pares ordenados
possíveis --- sem deduplicar A$\to$B e B$\to$A:

\begin{equation}
    \delta_{dir} = \frac{|E|}{|V|(|V|-1)}
    \label{eq:densidade}
\end{equation}

O denominador $|V|(|V|-1)$ representa exatamente o número de arcos dirigidos possíveis excluindo
self-loops. Isso difere da fórmula não-dirigida $2|E|/(|V|(|V|-1))$, que pressupõe pares
não-ordenados e subestimaria a conectividade real do grafo ao ignorar a assimetria das rotas.

% ─── 4. Métricas Globais ──────────────────────────────────────────────────────
\section{Métricas Globais e por Região}

\subsection{Métricas Globais}

\begin{table}[H]
\centering
\caption{Métricas globais do grafo}
\begin{tabular}{ll}
\toprule
\textbf{Métrica} & \textbf{Valor} \\
\midrule
Número de nós ($|V|$)               & 144 aeroportos \\
Número de arestas dirigidas ($|E|$)  & 955 arcos \\
Densidade global ($\delta_{dir}$)    & 0{,}0464 \\
\bottomrule
\end{tabular}
\end{table}

\subsection{Métricas por Região}

\begin{table}[H]
\centering
\caption{Métricas por região geográfica}
\begin{tabular}{lccc}
\toprule
\textbf{Região} & \textbf{Nós ($|V|$)} & \textbf{Arestas ($|E|$)} & \textbf{Densidade} \\
\midrule
Sudeste      & 38 & 162 & 0{,}1152 \\
Nordeste     & 32 &  89 & 0{,}0897 \\
Centro-Oeste & 15 &  17 & 0{,}0810 \\
Sul          & 23 &  36 & 0{,}0711 \\
Norte        & 36 &  77 & 0{,}0611 \\
\bottomrule
\end{tabular}
\end{table}

A região Sudeste apresenta a maior densidade (0{,}1152), refletindo a concentração de hubs
(\texttt{SBGR}, \texttt{SBKP}, \texttt{SBGL}, \texttt{SBSP}) e a alta frequência de voos
intraregionais. O Norte, apesar de ter o maior número de aeroportos (36), apresenta a menor
densidade (0{,}0611) --- consequência da dispersão geográfica e da baixa frequência de conexões
diretas entre cidades da região, que dependem fortemente de hubs como Manaus (\texttt{SBEG}) e
Belém (\texttt{SBBE}) para interligar destinos mais distantes. O Centro-Oeste, menor em número
de aeroportos (15), apresenta densidade de 0{,}0810 --- superior à do Norte e Sul, evidenciando
que seus aeroportos estão relativamente bem interconectados dentro da região.

\subsection{Rankings}

\begin{itemize}
    \item \textbf{Aeroporto mais conectado:} \texttt{SBKP} (Campinas/Viracopos) --- grau 63.
    \item \textbf{Maior fluxo inter-regional:} Sudeste $\rightarrow$ Sudeste com 14.214 voos em
          janeiro de 2026.
\end{itemize}

% ─── 5. Visualizações Analíticas ─────────────────────────────────────────────
\section{Visualizações Analíticas}

\subsection{Distribuição de Graus}

\begin{figure}[H]
    \centering
    \includegraphics[width=0.85\textwidth]{out/distribuicao_graus.png}
    \caption{Distribuição de graus dos aeroportos}
    \label{fig:dist_graus}
\end{figure}

A distribuição de graus revela uma malha com poucos hubs altamente conectados e muitos aeroportos
periféricos de baixo grau --- padrão típico de redes de escala livre (\emph{scale-free}). A maioria
dos aeroportos opera apenas algumas rotas, enquanto Campinas (SBKP), Guarulhos (SBGR) e outros
grandes hubs concentram a maior parte das conexões.

\subsection{Ranking dos Mais Conectados}

\begin{figure}[H]
    \centering
    \includegraphics[width=0.85\textwidth]{out/ranking_graus.png}
    \caption{Ranking dos aeroportos por grau}
    \label{fig:ranking}
\end{figure}

\subsection{Densidade por Região}

\begin{figure}[H]
    \centering
    \includegraphics[width=0.75\textwidth]{out/densidade_regioes.png}
    \caption{Densidade da sub-rede por região geográfica}
    \label{fig:densidade}
\end{figure}

\subsection{Heatmap de Fluxo Inter-Regional}

\begin{figure}[H]
    \centering
    \includegraphics[width=0.75\textwidth]{out/fluxo_regioes.png}
    \caption{Heatmap de fluxo de voos entre pares de regiões em janeiro/2026}
    \label{fig:heatmap}
\end{figure}

O heatmap evidencia o eixo Sudeste--Nordeste como o corredor inter-regional mais movimentado
(9.382 voos), seguido de Sudeste--Sul (6.935 voos). Fluxos envolvendo o Norte são os menores
entre regiões distintas.

% ─── 6. Caminhos Mínimos — Dijkstra ─────────────────────────────────────────
\section{Caminhos Mínimos --- Dijkstra}

O algoritmo de Dijkstra foi executado sobre os 5 pares de origem-destino exigidos pelo enunciado.
O custo reportado é a soma dos pesos $w(u,v)$ ao longo do caminho encontrado.

\begin{table}[H]
\centering
\caption{Resultados do Dijkstra para os pares obrigatórios e adicionais}
\begin{tabular}{llcl}
\toprule
\textbf{Origem} & \textbf{Destino} & \textbf{Custo} & \textbf{Caminho} \\
\midrule
SBEG & SBGR & --- & --- \\
SBRF & SBPA & --- & --- \\
SBPV & SBVT & --- & --- \\
SBRB & SBJP & --- & --- \\
SBGO & SBBE & --- & --- \\
\bottomrule
\end{tabular}
\end{table}

\noindent\textit{(Preencher com os valores gerados por \texttt{make solve} antes da entrega.)}

% ─── 7. Decisões de Implementação ────────────────────────────────────────────
\section{Decisões de Implementação}

\subsection{Identificadores: ICAO em vez de IATA}

Utilizamos códigos \textbf{ICAO} (ex.: \texttt{SBGR}) e não IATA (ex.: \texttt{GRU}) porque a
biblioteca \texttt{airportsdata} usa ICAO como chave primária e o dataset VRA já registra os
aeroportos nesse padrão.

\subsection{Grafo Interativo}

Optamos por não gerar o arquivo \texttt{out/grafo\_interativo.html} via pyvis. Em seu lugar,
desenvolvemos um frontend em \textbf{Next.js} com mapa SVG, projeção Mercator implementada
manualmente, painéis interativos de BFS, DFS e Dijkstra, filtro por região e bottom sheet com
gráficos analíticos (Recharts). A experiência entregue é substancialmente mais rica do que
qualquer output estático de pyvis permitiria.

\subsection{Fórmula de Densidade}

Utilizamos a fórmula dirigida (Equação~\ref{eq:densidade}), consistente com a natureza do grafo.
Essa fórmula é aplicada uniformemente nas métricas globais, regionais e de ego-rede, tanto no
backend Python quanto no frontend TypeScript.

% ─── 8. Conclusões ───────────────────────────────────────────────────────────
\section{Conclusões}

A modelagem da malha aérea brasileira como grafo dirigido ponderado por frequência de voos captura
de forma fiel a assimetria das rotas e a intensidade operacional de cada conexão. O uso do
logaritmo natural na fórmula de peso suaviza a escala sem perder a ordenação relativa entre
conexões raras e corredores de alta demanda.

A análise revela uma rede com características de escala livre: poucos hubs (SBKP, SBGR, SBGL)
concentram grau desproporcional, enquanto a maioria dos 144 aeroportos opera poucas rotas. O
corredor Sudeste--Nordeste domina o fluxo inter-regional, e o Norte apresenta a menor densidade
de sub-rede, refletindo as particularidades geográficas da região.

\end{document}
